\chapter{PIC architecture}
\label{ch:picarch}

\section{Block boundary}

\file{hdl/pic.v} is a self-contained IP block. Its interface is 16 hardware
interrupt lines, the four-signal CPU handshake, and one AXI4-Lite slave port for
configuration and status. It holds no CPU state and makes no assumption about the
CPU beyond the claim/end-of-interrupt protocol.

The feature specification for this block was explicitly a draft that delegated
the ``non-standard design challenges'' to the implementer. The choices recorded
in Chapter~\ref{ch:picbehav} therefore \emph{are} the feature specification for
those areas.

\begin{table}[H]
\centering
\caption{PIC top-level interface (\file{hdl/pic.v})}
\label{tab:picif}
\begin{tabular}{@{}llL{0.48\linewidth}@{}}
\toprule
\textbf{Signal} & \textbf{Width} & \textbf{Function} \\
\midrule
\sig{irq\_src}      & 16 & Hardware interrupt lines, level- or edge-triggered per source. \\
\sig{cpu\_irq}      & 1  & Registered level request to the CPU. \\
\sig{cpu\_irq\_vec} & 4  & Identifier of the source being offered. \\
\sig{cpu\_irq\_ack} & 1  & Claim pulse from the CPU. \\
\sig{cpu\_irq\_eoi} & 1  & End-of-interrupt pulse from the CPU. \\
\sig{s\_axi\_*}     & -- & AXI4-Lite slave, all five channels. \\
\bottomrule
\end{tabular}
\end{table}

Two internal sizes are fixed by localparam: \sig{NSRC} = 16 logical source slots,
and \sig{MAXNEST} = 16, the physical depth of the nesting stack.

\section{Source model}

Each of the 16 slots carries a hardware line and a software channel, giving 32
logical sources feeding one resolution pipeline. The slot request is
(\file{hdl/pic.v}):

\begin{center}
\ttfamily\small
hw\_req = cfg\_trig[s] ? edge\_pend[s] : irq\_src[s];\\
req[s]  = (hw\_req | sw\_pend[s]) \& int\_enable[s];
\end{center}

A software request merges into the same slot, so once inside the pipeline it is
indistinguishable from a hardware one --- same priority, same deadline, same
nesting behaviour.

\section{Resolution pipeline}

\begin{figure}[H]
\centering
\begin{tikzpicture}[node distance=6mm]
  \node[blkf,minimum width=30mm,align=left] (slot)
    {\textbf{per slot, $\times$16}\\[2pt]
     \scriptsize level / edge detect\\
     \scriptsize $\vert$ software channel\\
     \scriptsize \& INT\_ENABLE\\[3pt]
     \scriptsize deadline counter\\
     \scriptsize $\to$ effective band};

  \node[blk,right=16mm of slot,minimum width=34mm,align=left] (key)
    {\textbf{priority key}\\[2pt]
     \scriptsize \{band\_urgency[1:0],\\
     \scriptsize \ \ INTRA[3:0], 15$-$index\}\\
     \scriptsize 10 bits, higher = more urgent};

  \node[blk,right=16mm of key,minimum width=32mm,align=left] (res)
    {\textbf{resolver}\\[2pt]
     \scriptsize eligible = req \& !active\\
     \scriptsize \& key $>$ top\_key\\
     \scriptsize masked if depth = NEST\_MAX};

  \node[blk,below=13mm of res,minimum width=32mm,align=left] (reg)
    {\textbf{output register}\\[2pt]
     \scriptsize cpu\_irq $\leftarrow$ offer\_val\\
     \scriptsize cpu\_irq\_vec $\leftarrow$ res\_id};

  \node[blkf,below=13mm of key,minimum width=34mm,align=left] (stack)
    {\textbf{nesting stack}\\[2pt]
     \scriptsize \{key, id\} pairs, depth 0..NEST\_MAX\\
     \scriptsize push on ack, pop on eoi};

  \node[blkf,below=13mm of slot,minimum width=30mm,align=left] (axi)
    {\textbf{axi\_lite\_slave}\\[2pt]
     \scriptsize register file\\
     \scriptsize 56 mapped words};

  \draw[flow] (slot) -- (key);
  \draw[flow] (key) -- (res);
  \draw[flow] (res) -- (reg);
  \draw[flow] (stack) -- node[lbl,right]{top\_key} (res);
  \draw[flow] (axi) -- node[lbl,above]{config} (slot);
  \draw[flow,<->] (axi) -- node[lbl,above]{status} (stack);
  \draw[flow] (reg.east) -- ++(9mm,0) node[lbl,right,align=left]{\scriptsize to CPU};
  \draw[flow] ($(stack.south)+(0,-1mm)$) ++(0,-4mm) -- (stack.south)
        node[lbl,below=5mm,align=center]{\scriptsize ack / eoi};
  \draw[flow] (slot.west) ++(-9mm,0) -- (slot.west) node[lbl,left=9mm]{\scriptsize irq\_src};
\end{tikzpicture}
\caption{PIC internal structure. The resolver is entirely combinational; the only
registered stage on the request path is the output pair.}
\label{fig:picblock}
\end{figure}

\subsection{Priority key}

Every source is reduced to a single 10-bit unsigned number and the resolver picks
the maximum:

\begin{center}
\ttfamily\small
key[s] = \{ band\_urg(eff\_band[s], band\_cfg), cfg\_intra[s], 15 - s \};
\end{center}

Band urgency dominates, then intra-band priority, then the last field. Because
that field is $15-\text{index}$, it is larger for smaller indices, so a tie is
won by the lowest-numbered source. The ordering is total: there is never an
ambiguous case and the winner is always unique.

Note what is \emph{not} in the key: the band \emph{number} carries no ranking of
its own. A band's urgency is looked up from \reg{BAND\_CONFIG}, which is what
lets software reorder all four bands with one register write.

\subsection{Eligibility and preemption}

\begin{center}
\ttfamily\small
if (req[i] \&\& !active[i] \&\& (!has\_active || key[i] > top\_key))
\end{center}

A source is eligible if it has a request, is not already in service, and --- when
anything is in service --- its key is \emph{strictly} greater than the key on top
of the nesting stack.

That single strict comparison does two jobs. It is the preemption rule: a more
urgent source displaces a less urgent handler. And it is the re-entrancy
guarantee: a source cannot re-offer itself, because its own key is not strictly
greater than itself.

\subsection{Output registration}

\sig{cpu\_irq} and \sig{cpu\_irq\_vec} are registered from the same resolver
output on the same clock edge, so the pair the CPU sees is always
self-consistent. The cost is one cycle of latency; the benefit is that the PIC
never sits on the combinational timing path of the source lines. \sig{cpu\_irq\_vec}
holds its last presented value while idle.

A second register, \sig{cpu\_irq\_vec\_d}, holds the vector as presented in the
previous cycle. Claims are attributed to \emph{that} copy, because
\sig{cpu\_irq\_ack} arrives one cycle after the CPU sampled the vector it acted
on. This keeps the PIC's notion of the claimed source aligned with the CPU's
latched \reg{mcause} even if a more urgent source appeared in the meantime.

\section{Per-source lifecycle}

The PIC has no state register per source. The lifecycle below is encoded in
separate flag arrays --- \sig{edge\_pend}, \sig{sw\_pend}, \sig{active},
\sig{escalated}, \sig{spurious} --- each with its own always block. The diagram
describes the behaviour those flags produce.

\begin{figure}[H]
\centering
\begin{tikzpicture}[x=1mm,y=1mm,>={Stealth[length=2.4mm]},semithick]
  \node[stst,minimum size=15] (idle) at (0,22)   {\scriptsize IDLE};
  \node[st,minimum size=18]   (pend) at (52,22)  {\scriptsize PENDING};
  \node[st,minimum size=17]   (act)  at (108,22) {\scriptsize ACTIVE};
  \node[st,minimum size=20]   (esc)  at (52,-18) {\scriptsize ESCALATED};

  % idle -> pending
  \draw[->,draw=accent] (idle) -- node[lbl,above,align=center]
    {line asserted, edge latched or\\\scriptsize software trigger, \&\ enabled} (pend);

  % pending -> active, two ways
  \draw[->,draw=accent] (pend) to[bend left=16]
    node[lbl,above,align=center,pos=0.5]{cpu\_irq\_ack\\\scriptsize (claim, push)} (act);
  \draw[->,draw=accent] (pend) to[bend right=16]
    node[lbl,below,align=center,pos=0.5]
    {request gone before the claim\\\scriptsize $\to$ SPURIOUS, still pushed} (act);

  % active -> idle, routed orthogonally above every other label
  \draw[->,draw=accent] (act.north) -- (108,52) -- (0,52) -- (idle.north);
  \node[lbl,above=0.5mm,align=center] at (54,52)
    {cpu\_irq\_eoi (pop) --- effective band and deadline counter\\
     \scriptsize reset to the configured values};

  % pending <-> escalated
  \draw[->,draw=accent] ($(pend.south)+(-5,0)$) -- ($(esc.north)+(-5,0)$)
    node[lbl,midway,left=1mm,align=right]{ddl\_cnt reaches\\\scriptsize DEADLINE};
  \draw[->,draw=accent] ($(esc.north)+(5,0)$) -- ($(pend.south)+(5,0)$)
    node[lbl,midway,right=1mm,align=left]{still pending,\\\scriptsize higher effective band};

  \draw[->,draw=accent] (esc) to[out=-130,in=-50,looseness=4.5]
    node[lbl,below=1mm,align=center]{MULTI = 1: counter re-arms, escalate again} (esc);
\end{tikzpicture}
\caption{Per-source lifecycle. Flag-encoded, not a state register: a source can be
pending and escalated simultaneously, and escalated condition only modifies the
effective band used to build the priority key.}
\label{fig:picsrcfsm}
\end{figure}

\section{Nesting stack}

The stack holds the \{key, id\} pair of every source currently in service:

\begin{center}
\ttfamily\small
reg [3:0] stack\_id  [0:MAXNEST-1];\\
reg [9:0] stack\_key [0:MAXNEST-1];\\
reg [4:0] depth;
\end{center}

\sig{depth} is the stack pointer and the only element that resets. The top is
\sig{top\_idx} = \sig{depth} $-$ 1, valid whenever \sig{depth} is non-zero.

Storing the key, not just the id, is what makes the preemption threshold stable:
reconfiguring a source's band while it is in service cannot corrupt nesting,
because the threshold in use is the snapshot taken at the claim. New pending
evaluations of that source use the new band immediately.

\begin{figure}[H]
\centering
\begin{tikzpicture}[>={Stealth[length=2.4mm]},semithick,node distance=26mm]
  \node[stst] (d0) {\scriptsize depth 0};
  \node[st,right=28mm of d0] (d1) {\scriptsize depth 1};
  \node[st,right=28mm of d1] (dn) {\scriptsize depth $n$};
  \node[st,right=28mm of dn,minimum size=16mm] (dmax) {\scriptsize NEST\_MAX};

  \draw[->,draw=accent] (d0) to[bend left=14] node[lbl,above]{ack: push} (d1);
  \draw[->,draw=accent] (d1) to[bend left=14] node[lbl,below]{eoi: pop} (d0);
  \draw[->,draw=accent,dashed] (d1) to[bend left=14] node[lbl,above]{ack} (dn);
  \draw[->,draw=accent,dashed] (dn) to[bend left=14] node[lbl,below]{eoi} (d1);
  \draw[->,draw=accent] (dn) to[bend left=14] node[lbl,above]{ack} (dmax);
  \draw[->,draw=accent] (dmax) to[bend left=14] node[lbl,below]{eoi} (dn);
  \draw[->,draw=accent,loop above] (dmax) to
     node[lbl,above,align=center]{offers masked,\\\scriptsize INT\_STATUS.OVF set} (dmax);
\end{tikzpicture}
\caption{Nesting depth. The limit is enforced by masking offers, so a claim that
would exceed it is never presented to the CPU in the first place.}
\label{fig:picdepth}
\end{figure}

\section{Register interface}

The AXI4-Lite handshake is not implemented in the PIC. It instantiates
\file{hdl/axi\_lite\_slave.v}, the shared slave used by both the PIC and the
machine timer (Section~\ref{sec:axislave}), and supplies only the register file:
a combinational read multiplexer, a one-cycle write pulse, and two
combinational qualifiers \sig{rd\_ok} / \sig{wr\_ok} that mark which offsets are
readable and writable. Everything else --- unmapped offsets, writes to read-only
registers --- becomes SLVERR through that module.
